\documentclass[10pt,a4paper]{extarticle}
\pagenumbering{gobble}

\usepackage{fontspec}
\setmainfont{TeX Gyre Termes}

\usepackage[turkish,shorthands=off]{babel}
\usepackage[left=0.6in,top=0.6in,right=0.6in,bottom=0.6in]{geometry}
\usepackage{hyperref}
\usepackage{titlesec}
\usepackage{enumitem}
\usepackage{datetime}
\usepackage{graphicx}
\titleformat{\section}{\large\bfseries\scshape\raggedright}{}{0em}{}[\titlerule]
\titlespacing*{\section}{0pt}{12pt}{8pt}
\begin{document}
\begin{center}
    \begin{minipage}{\textwidth}
        \centering
        {\LARGE\textbf{Yiğit Efe Avcı}} \hspace{2pt} {\LARGE{Full-Stack Geliştirici}}\\[10pt]
        \href{mailto:yigit433.devs@gmail.com}{yigit433.devs@gmail.com} \textbullet\
        \href{https://github.com/yigit433}{github.com/yigit433} \textbullet\
        \href{https://linkedin.com/in/yigitefeavci}{linkedin.com/in/yigitefeavci} \textbullet\
        \href{https://yigit433.vercel.app}{yigit433.vercel.app}
    \end{minipage}
\end{center}
\section{Hakkımda}
Veri bilimi, yazılım geliştirme ve istatistiksel modelleme alanlarında güçlü bir ilgiye sahip istatistik lisans öğrencisiyim. Python, R ve SPSS gibi araçları kullanarak verileri eyleme dönüştürülebilir içgörülere dönüştürme konusunda uzmanlaşıyorum. Teknik odağım, React, Next.js, Go ve PostgreSQL gibi teknolojileri kullanarak ölçeklenebilir, verimli ve kullanıcı odaklı web uygulamaları oluşturduğum full-stack geliştirmeye uzanıyor.
\section{Deneyim}
\textbf{\href{https://codeshare.me}{Code Share}} \hfill Londra, Birleşik Krallık\\
\textit{Yönetici Ortak (Serbest Meslek)} \hfill Ağustos 2020 -- Mayıs 2025\\
\textit{Yazılım Geliştirici (Freelance, Uzaktan)} \hfill Haziran 2019 -- Haziran 2023\\[6pt]
\textbf{WiFiCeps} \hfill Türkiye\\
\textit{Backend Developer (Proje Bazlı)} \hfill Ocak 2023 -- Şubat 2023
\begin{itemize}[leftmargin=*,noitemsep,topsep=0pt]
    \item WiFiCeps eğitim platformu için özel bir medya sunucusu geliştirdim.
    \item Google Drive'a yüklenen videoları izleyen, güncellemeleri tespit edip platformdaki mevcut sürümlerle otomatik olarak senkronize eden bir sistem tasarladım.
    \item Google Drive'dan silinen videoları medya sunucusundan otomatik olarak kaldıran bir süreç uyguladım.
    \item Python tabanlı bir otomasyon çerçevesi ve FFmpeg destekli bir video işleme algoritması inşa ettim.
    \item Otomatik işleme süreci ile video boyut ve kalitesini optimize ederek dosyaları yaklaşık \textasciitilde50MB seviyesine indirdim.
    \item Oynatma sorunlarını belirgin şekilde azalttım ve genel akış performansını iyileştirdim.
\end{itemize}
\textbf{Konda Araştırma ve Danışmanlık} \hfill Türkiye\\
\textit{Saha Araştırmacısı (Anketör)} \hfill Ay Yıl -- Ay Yıl
\begin{itemize}[leftmargin=*,noitemsep,topsep=0pt]
    \item Yurt çapında yürütülen anketler için saha çalışmaları gerçekleştirdim; yüz yüze görüşmeler ve veri toplama süreçlerini kapsıyordu.
    \item Metodolojik standartlara uyarak ve katılımcılarla doğrudan etkileşim kurarak araştırma projelerine doğru ve güvenilir veri girdisi sağladım.
\end{itemize}
\section{Projeler}
\textbf{\href{https://pickin.co}{Pickin} — Gerçek zamanlı çekiliş platformu} \hfill Next.js, TypeScript, PostgreSQL\\
\begin{itemize}[leftmargin=*,noitemsep,topsep=0pt]
    \item Redis Pub/Sub ile beslenen Server-Sent Events üzerinden anlık katılımcı ve durum güncellemeleri.
    \item NATS JetStream iş kuyruğu ile zamanlanmış çekiliş başlatma/bitirme ve otomatik kazanan seçimi.
    \item OAuth tabanlı katılım şartları (YouTube abonelik/beğeni, Twitch takip) — Better Auth ile entegre.
    \item RBAC sistemi (4 rol, 15 izin), çoklu dil (TR/EN), QR paylaşım, OBS yayın görünümü, CSV dışa aktarım.
    \item Stack: Next.js 16 (App Router, Turbopack), Drizzle ORM, Tailwind CSS v4, Framer Motion, TipTap, Docker.
\end{itemize}
\textbf{\href{https://deckplane.com}{Deckplane} — Dağıtık Docker yönetim platformu} \hfill Bun, Hono, TypeScript, PostgreSQL\\
\begin{itemize}[leftmargin=*,noitemsep,topsep=0pt]
    \item Control-plane / host-agent mimarisi: her Docker host'unda çalışan ajanlar, merkezi API ile WebSocket ve REST üzerinden haberleşir.
    \item Birden fazla sunucudaki container, image, network ve volume'leri tek API üzerinden ve asenkron komut kuyruğu ile yönetir.
    \item WebSocket tabanlı interaktif container terminali (oturum kaydı dahil), gerçek zamanlı log ve CPU/bellek istatistik akışı.
    \item GitHub repolarından veya docker-compose arşivlerinden dağıtım otomasyonu; OAuth ile GitHub entegrasyonu (branch, secret, webhook).
    \item RBAC (4 rol), çoklu kiracılık, değişmez denetim kayıtları, token-bucket rate limiting, otomatik OpenAPI 3.0 + Swagger UI.
    \item Stack: Bun, Hono, Drizzle ORM, JWT (jose) + Argon2, Zod, Pino, AES-256-GCM, Traefik (HTTPS + WebSocket).
\end{itemize}
\textbf{\href{https://github.com/yigit433/kommando}{kommando} — Go için CLI araç kiti} \hfill Go\\
\begin{itemize}[leftmargin=*,noitemsep,topsep=0pt]
    \item Alt komut ve flag desteği ile komut satırı uygulamaları geliştirmek için hafif bir kütüphane.
    \item Ergonomik API tasarımına ve dış bağımlılık olmamasına odaklı.
\end{itemize}
\textbf{\href{https://github.com/yigit433/helper-tui}{helper-tui} — Terminal arayüzü yardımcısı} \hfill Rust\\
\begin{itemize}[leftmargin=*,noitemsep,topsep=0pt]
    \item Geliştiricilerin sık yaptığı işlemleri terminalden hızlandıran çoklu platform TUI uygulaması.
    \item Rust ile, hızlı tepki ve minimum ikili boyut hedeflenerek yazıldı.
\end{itemize}
\textbf{\href{https://github.com/yigit433/feather-stream}{feather-stream} — Google Drive medya oynatıcı} \hfill Tauri, React, TypeScript\\
\begin{itemize}[leftmargin=*,noitemsep,topsep=0pt]
    \item Google Drive'da depolanan medya dosyalarını izlemek için masaüstü uygulaması.
    \item Tauri (Rust çekirdek) ile React/TypeScript arayüzü kullanılarak uçtan uca geliştirildi.
\end{itemize}
\textbf{\href{https://github.com/yigit433/jsonparser}{jsonparser} — C ile yazılmış JSON parser} \hfill C\\
\begin{itemize}[leftmargin=*,noitemsep,topsep=0pt]
    \item Sıfırdan, C dilinde yazılmış JSON belge ayrıştırıcısı ve işleyicisi.
    \item Elle yazılmış tokenizer ve özyinelemeli iniş ayrıştırıcı; dış bağımlılık yok.
\end{itemize}
\section{Sertifikalar}
\begin{itemize}[leftmargin=*,noitemsep,topsep=0pt]
    \item \textbf{Claude Code in Action} --- Anthropic, Nisan 2026. Credential ID: qenfi6q2dx7k.
    \item \textbf{AI Agent Geliştirme ve MCP Temelleri} --- Siber Kulüpler Birliği, Ocak 2026. Credential ID: SKB-2026-733OF.
    \item \textbf{Sıfırdan İleri Seviye Python Programlama} --- BTK Akademi, Ocak 2025. Credential ID: OKMhqXvw04.
    \item \textbf{Google Flutter ile Mobil Uygulama} --- BTK Akademi, Ocak 2025. Credential ID: L8dcn9b8V6.
    \item \textbf{Machine Learning with Python} --- IBM, Şubat 2025. Credential ID: 1H9MA1TRY556.
\end{itemize}
\section{Yetenekler}
\textbf{Diller:} Python, R, TypeScript, JavaScript, GoLang, Rust, Dart \\
\textbf{Veri \& ML:} Pandas, Polars, NumPy, Matplotlib, Seaborn, Scikit-learn, TensorFlow, PyTorch, SPSS \\
\textbf{Backend:} Bun, Node.js, Express.js, REST API \\
\textbf{Frontend:} React, Next.js, Tailwind CSS, Framer Motion \\
\textbf{Mobil / Çoklu Platform:} Flutter, React Native, Tauri \\
\textbf{Veritabanları:} PostgreSQL, MongoDB, MySQL, Supabase \\
\textbf{Araçlar:} Git, Docker, Drizzle ORM, Postman
\section{Eğitim}
\textbf{İstanbul Ticaret Üniversitesi} \hfill İstanbul, Türkiye\\
\textit{İstatistik Lisans} \hfill 2023 -- 2027\\
\end{document}
